% 中文摘要
\chapter*{\hfill 中 文 摘 要 \hfill}
\addcontentsline{toc}{chapter}{中文摘要}
\thispagestyle{fancy}

\setcounter{page}{1}
\pagenumbering{Roman}

\par SCADA（监控与数据采集）系统作为国家关键基础设施的神经中枢，广泛应用于燃气、
电力、水务等工业控制领域，其安全性直接关系到公共安全与经济稳定。随着工业互联网
与人工智能技术的深度融合，传统基于规则与签名的入侵检测方法已难以应对日益复杂的
网络攻击，亟需研究面向 SCADA 系统的智能化、轻量化入侵检测方法。

\par 然而，SCADA 工控场景对入侵检测模型提出了独特而严苛的要求：一方面，工业控
制协议（如 Modbus）的时序结构蕴含丰富的领域语义，通用时序模型难以充分利用；
另一方面，受限于边缘设备的算力与存储，检测模型必须兼顾精度与部署可行性。现有
研究多沿用通用网络入侵检测的特征体系，缺乏协议层语义编码，且鲜有工作系统性地
探索从算法设计到芯片部署的完整优化链条。

\par 针对上述问题，本文以 IanArffDataset 数据集（274{,}000 条 SCADA 燃气管道记
录，20 维原始特征，三层标签）为研究对象，围绕``轻量化入侵检测''这一核心目标，
开展了从数据特征、模型架构到边缘部署的全栈研究，主要贡献与创新点如下：

\par （1）提出了一种基于 SCADA 协议语义的领域特征工程方法。深入分析 Modbus 协
议的 4 行主-从交互结构，将 4 行视为一个 Modbus 周期，在 17 维最小特征的基础上
设计 27 维协议语义特征（3 行级 + 8 窗口级）。在 LightGBM、TCN 等多个模型上的对
比实验表明，所提特征集在所有模型上一致优于原始 44 维特征，验证了``领域特征工
程的边际收益高于模型升级''的结论。

\par （2）提出了一种融合 SE 注意力的 TCN 轻量入侵检测模型（TCN-SE）。在系统
对比 27 种典型时序模型的基础上，发现 TCN 因果卷积的归纳偏置天然适配 Modbus
短时序结构；进一步引入 SE 通道注意力，仅增加 3{,}300 参数（+4.3\%）即在测试
集上取得 Macro-F1 = 0.8340，PR-AUC = 0.9025，首次在该任务上实现深度学习模型对
梯度提升树（LightGBM，Macro-F1 = 0.8337）的超越。通过窗口长度（$w=4, 8, 16$）
与通道宽度（ch=8 $\sim$ 64）的消融实验，揭示了``8 步窗口（2 个 Modbus 周期）
是该数据集的 Pareto 最优结构''，并验证了 SE 模块对少数类召回率的提升作用。

\par （3）提出了一种面向 MCU 部署的极限压缩与混合精度量化方案。通过系统的通道
宽度消融，得到参数量仅 4{,}237 的 TCN V19 变体，PR-AUC 反升至 0.9133（+0.67\%），
挑战了``模型越大性能越好''的传统认知。进一步提出``权重 INT8 + 激活 FP32''的
混合精度量化方案，在 Cortex-M4 168\,MHz 平台上实现 5.5\,KB 模型、0.47\,ms 推理
延迟，量化前后预测一致率达 99.80\%，SNR $\geq$ 38\,dB 无损。完整 380 行 C99 推理
代码已开源，可直接部署于 STM32、ESP32、Arduino 等嵌入式平台。

\par 本文围绕``数据-算法-芯片''三层递进的研究路线，形成了从领域特征到 MCU 部署
的完整创新链条。三个创新点在公开数据集上取得了 Macro-F1 = 0.8340（单模）与
PR-AUC = 0.9133（压缩单模）的最优结果，并验证了在 5.5\,KB 嵌入式模型上保持工
业级检测性能的可行性，为 SCADA 工控入侵检测的工程化落地提供了可复现的参考方
案。

\vspace{1em}

\noindent\textbf{关键词：} SCADA；入侵检测；时序卷积网络；模型压缩；混合精度量化；嵌入式部署

\clearpage
